\chapter{Summary and outlook}
\label{ch:conclusion}
In this study, we present a theoretical investigation of excitonic transitions in two--dimensional transition--metal dichalcogenides under photoexcitation by scalar and vector vortex beams.

In Chapter 2, we show that the well-defined orbital angular momentum (OAM) of $\ell \hslash$ carried by Laguerre--Gaussian (LG) beams is preserved under the paraxial approximation. Upon transformation to the Coulomb gauge, coupling with spin angular momentum (SAM) $\sigma \hslash$ naturally emerges. The adoption of the Coulomb gauge is justified within the framework of light--matter interaction, and the resulting spin--orbit interaction (SOI) term, $\sigma+\ell$, appears in the phase of the longitudinal component of LG beams in the angular spectrum representation. Furthermore, by superposing two LG beams to form vector vortex beams (VVBs), the SOI term is preserved as $\Delta J=(\sigma'+\ell')-(\sigma+\ell)$ in the amplitude of the longitudinal field. This quantity governs the variation of the relative phase along the azimuthal angle. Thus, when $\Delta J=0$, the azimuthal structure vanishes, and the interference is fully controlled by the relative phase between the two LG beams, leading to either complete constructive or destructive outcomes.

In Chapter 3, Fermi's golden rule is employed to evaluate the transition rates of excitons, focusing on valley-mixed bright excitons (BXs) and gray excitons (GXs) under VVB excitation. Due to the approximately isotropic momentum-space distribution and out-of-plane orientation of the GX transition dipole, the transition rate directly reflects the longitudinal component of the VVB. As a result, selective excitation of GXs is achieved under the condition $\Delta J=0$, corresponding to complete switching of the longitudinal field. On the other hand, valley-mixed BXs exhibits behavior similar to GXs. However, the condition $\Delta J=0$ corresponds to matching between the BX transition dipole and the light polarization. Although selective excitation of valley-mixed BXs is theoretically possible, the few-meV valley splitting poses challenges for experimental verification.

In Chapter 4, first--order time--dependent perturbation theory is applied to construct a superposition of finite-momentum exciton (SFME) states, with the aim of determining the static envelope function of the exciton wave packet and examining the resulting spatial patterns. It is found that a global phase in valley-mixed BXs modifies the resulting envelope function, necessitating an explicit treatment of the center--of--mass momentum ($\boldsymbol{Q}$)--dependent phase structure of the exciton wave function. The extraction, based on zeroth-order and related approximations, is successful for valley excitons. In contrast, valley-mixed BXs introduce additional phase complexity originating from their internal structure. Simulations for valley excitons further demonstrate that phase factors play a crucial role in shaping the static envelope function.

Future work may extend this framework to other VVB configurations and incorporate full exciton wave functions. The fascinating interplay between structured light and matter enables further control of exciton dynamics, with potential applications in valleytronics and optoelectronic systems.
